% 第1章 引言（Introduction）
\chapter{引言}

本书讨论的是磁性在凝聚态（condensed matter）中的种种表现。固体中含有大量磁矩
（magnetic moments），它们能够协同行动，产生与彼此孤立时截然不同的行为；再加上
可以存在的磁相互作用类型多种多样，这就使得真实体系的磁性质丰富得惊人。本书的
思路是放缓节奏、逐块搭建这幅图景。在这一导论性的章节里，我们先回顾初等经典物理
和量子物理中关于磁矩的一些事实；随后各章依次讨论：第二章讨论当大量磁矩被置入
固体但彼此隔离、也与环境隔离时的行为；第三章考虑它们的近邻环境所产生的影响；
接着第四章讨论磁矩之间可能存在的各种磁相互作用。到第五章我们就可以讨论长程序
（long range order）的出现，第六章讨论它与对称性破缺（broken symmetry）概念的
联系。最后两章把这一概念的含义贯彻到各种不同情形之中。全书采用 SI 单位制
（cgs 单位制的描述与换算表见附录 A）。

\section{磁矩}

磁学中最基本的对象是磁矩。在经典电磁学中，可以把磁矩等同于一个电流环。若有电流
$I$ 绕一个元（即无穷小的）有向环路流动，环面积为 $|d\mathbf{S}|$（见图~1.1(a)），
则磁矩 $d\boldsymbol{\mu}$ 为
\begin{equation*}
\mathbf{d}\boldsymbol{\mu} = I\,\mathbf{dS},\tag{1.1}
\end{equation*}
磁矩的单位是 A\,m$^2$。矢量 $\mathbf{dS}$ 的长度等于环的面积，其方向垂直于环面，
指向由环绕元电流的方向按右手定则确定。

这一对象也等价于一个磁偶极子（magnetic dipole）；之所以这样称呼，是因为它的行为
类似于电偶极子（相隔一小段距离的一正一负两个电荷）。因此可以设想：磁偶极子是这样
一种对象——两个带相反磁荷的磁单极子（magnetic monopoles），沿 $\mathbf{dS}$ 方向
相隔一小段距离（关于电磁学的背景知识见附录 B）。

\begin{figure}
\centering
\includegraphics[width=0.62\textwidth]{fig1_01.pdf}
\caption{（a）元磁矩：由元电流环产生的 $\mathrm{d}\boldsymbol{\mu}=I\,\mathrm{d}\mathbf{S}$。（b）一个有限回路（俯视图）上的电流 $I$ 的磁矩 $\boldsymbol{\mu}=I\int\mathrm{d}\mathbf{S}$，可由把许多无穷小电流环的磁矩求和得到。}
\end{figure}

磁矩 $d\boldsymbol{\mu}$ 垂直于电流环平面，因而既可以与环上电荷的角动量矢量平行，
也可以反平行。对于有限大小的回路，可以把其磁矩 $\boldsymbol{\mu}$ 表示为分布在
回路面积上的许多相同的无穷小电流环的磁矩之和（见图~1.1(b)）：相邻无穷小环的电流
彼此抵消，只留下沿回路周界流动的电流。于是
\begin{equation*}
\boldsymbol{\mu} = \int \mathrm{d}\boldsymbol{\mu} = I\int \mathrm{d}\mathbf{S}.\tag{1.2}
\end{equation*}

\subsection{磁矩与角动量}

电流环源于电荷的运动。本书考虑的所有电荷都与具有质量的粒子相联系，因此书中所有
电流环里质量也在作轨道运动——磁矩总是与角动量相伴随。

在原子中，沿轨道运动的电子所具有的磁矩 $\boldsymbol{\mu}$ 与该电子的角动量
$\mathbf{L}$ 方向一致，且与它成正比。我们写成
\begin{equation*}
\boldsymbol{\mu} = \gamma\,\mathbf{L},\tag{1.3}
\end{equation*}
其中 $\gamma$ 是一个常数，称为旋磁比（gyromagnetic ratio）。磁矩与角动量之间的
这一关系由 1915 年发现的爱因斯坦--德哈斯效应（Einstein--de Haas effect）所证实：
一根铁磁棒沿其轴向用细丝竖直悬挂（见图~1.2），起初静止且未磁化，随后沿棒长方向
加竖直磁场使其磁化。这一竖直磁化来自原子磁矩的排列，对应一定的角动量；为守恒
总角动量，铁棒将开始沿相反方向绕轴转动。测出棒的角动量，便可推知原子磁矩所携带
的角动量，从而得到旋磁比。爱因斯坦--德哈斯效应是磁化诱导的转动；还存在逆效应，
即巴尼特效应（Barnett effect）——转动诱导磁化。两种现象都表明磁矩与角动量相联系。
\bio{Albert Einstein（1879--1955）\\Wander J. de Haas（1878--1960）}

\begin{figure}
\centering
\includegraphics[width=0.5\textwidth]{fig1_02.pdf}
\caption{爱因斯坦--德哈斯效应。铁磁棒悬挂于细丝之下。线圈提供磁场使铁磁体磁化并产生
转动。实验可以共振方式实现：周期性地反转线圈中的电流（从而反转铁磁体中的磁化），
观察角响应随频率的变化。}
\end{figure}

\subsection{进动}

现在考虑处于磁场 $\mathbf{B}$ 中的磁矩 $\boldsymbol{\mu}$（见图~1.3）。磁矩的能量为
\begin{equation*}
E = -\boldsymbol{\mu}\cdot\mathbf{B},\tag{1.4}
\end{equation*}
（见附录 B）因此磁矩沿磁场方向时能量最低。磁矩还会受到力矩（torque）
\begin{equation*}
\mathbf{G} = \boldsymbol{\mu}\times\mathbf{B},\tag{1.5}
\end{equation*}
（见附录 B）——假如磁矩不携带任何角动量，这个力矩将把磁矩转向磁场方向\fn{1}{对于处于电场 $\boldsymbol{\mathcal{E}}$ 中的电偶极子 $\mathbf{p}$，能量为 $E=-\mathbf{p}\cdot\boldsymbol{\mathcal{E}}$，力矩为 $\mathbf{G}=\mathbf{p}\times\boldsymbol{\mathcal{E}}$。静止的电偶极矩只是两个相隔一定距离的静止电荷，不携带任何角动量，因此若 $\boldsymbol{\mathcal{E}}$ 与 $\mathbf{p}$ 不共线，力矩 $\mathbf{G}$ 将把 $\mathbf{p}$ 转向 $\boldsymbol{\mathcal{E}}$。而静止的磁矩则与角动量相联系，故行为不同。}。

然而，由于磁矩通过式~(1.3) 与角动量 $\mathbf{L}$ 相联系，而力矩等于角动量的时间
变化率，式~(1.5) 可改写为
\begin{equation*}
\frac{\mathrm{d}\boldsymbol{\mu}}{\mathrm{d}t} = \gamma\,\boldsymbol{\mu}\times\mathbf{B}.\tag{1.6}
\end{equation*}
这意味着 $\boldsymbol{\mu}$ 的变化同时垂直于 $\boldsymbol{\mu}$ 和 $\mathbf{B}$。
磁场并不把 $\boldsymbol{\mu}$ 拉向 $\mathbf{B}$，而是使 $\boldsymbol{\mu}$ 的方向
绕 $\mathbf{B}$ 进动（precession）。式~(1.6) 还表明 $|\boldsymbol{\mu}|$ 不随时间
改变。这一情形与陀螺的旋转完全类似。\fn{2}{设想一个自转轴相对竖直方向倾斜的陀螺。向下的重力对陀螺施加一个（水平的）力矩。若陀螺不自转，它只会倒下；但正因为它在自转，它具有平行于自转轴的角动量，力矩使自转轴在水平面内平行于力矩方向运动——陀螺进动了。}
\bio{Joseph Larmor（1857--1942）}

下面的例子将针对一个具体情形详细求解式~(1.6)。

\begin{example}{1.1}
设 $\mathbf{B}$ 沿 $z$ 方向，$\boldsymbol{\mu}$ 初始时与 $\mathbf{B}$ 成 $\theta$ 角
并位于 $xz$ 平面内（见图~1.4）。则
\begin{align*}
\dot{\mu}_x &= \gamma B\mu_y\tag{1.7}\\
\dot{\mu}_y &= -\gamma B\mu_x\tag{1.8}\\
\dot{\mu}_z &= 0,\tag{1.9}
\end{align*}
所以 $\mu_z$ 不随时间变化，而 $\mu_x$、$\mu_y$ 都在振荡。解这组微分方程得
\begin{align*}
\mu_x(t) &= |\boldsymbol{\mu}|\sin\theta\cos(\omega_{\mathrm{L}}t)\tag{1.10}\\
\mu_y(t) &= |\boldsymbol{\mu}|\sin\theta\sin(\omega_{\mathrm{L}}t)\tag{1.11}\\
\mu_z(t) &= |\boldsymbol{\mu}|\cos\theta,\tag{1.12}
\end{align*}
其中
\begin{equation*}
\omega_{\mathrm{L}} = \gamma B\tag{1.13}
\end{equation*}
称为拉莫尔进动频率（Larmor precession frequency）。
\end{example}

\begin{figure}
\centering
\includegraphics[width=0.32\textwidth]{fig1_03.pdf}
\caption{处于磁场 $\mathbf{B}$ 中的磁矩 $\boldsymbol{\mu}$，能量为
$-\boldsymbol{\mu}\cdot\mathbf{B}=-\mu B\cos\theta$。}
\end{figure}

\begin{figure}
\centering
\includegraphics[width=0.4\textwidth]{fig1_04.pdf}
\caption{磁场 $\mathbf{B}$ 中的磁矩 $\boldsymbol{\mu}$ 以拉莫尔进动频率 $\gamma B$ 绕磁场
进动，$\gamma$ 为旋磁比。$\mathbf{B}$ 沿 $z$ 轴，磁矩初始时位于 $xz$ 平面内、与
$\mathbf{B}$ 成 $\theta$ 角。磁矩绕半角为 $\theta$ 的圆锥进动。}
\end{figure}

注意旋磁比 $\gamma$ 是把角动量与磁矩（通过式~1.3）、又把进动频率与磁场（通过式~1.13）
联系起来的比例常数。进动现象暗示了后面内容的精微之处：磁场不仅使磁矩排列整齐，
还能诱发各种动力学效应。

\subsection{玻尔磁子}

在继续之前，值得快速估算一下原子磁矩的大小，从而推知旋磁比的量级。考虑一个电子
（电荷 $-e$，质量 $m_{\mathrm{e}}$）绕氢原子核作圆周轨道运动（见图~1.5）。绕核
电流为 $I=-e/\tau$，其中 $\tau=2\pi r/v$ 是轨道周期，$v=|\mathbf{v}|$ 是速率，
$r$ 是圆轨道半径。基态中电子角动量的大小 $m_{\mathrm{e}}vr$ 必等于 $\hbar$，于是
电子的磁矩为
\begin{equation*}
\mu = \pi r^{2} I = -\frac{e\hbar}{2m_{\mathrm{e}}} \equiv -\mu_{\mathrm{B}}\tag{1.14}
\end{equation*}
其中 $\mu_{\mathrm{B}}$ 是玻尔磁子（Bohr magneton），
\begin{equation*}
\mu_{\mathrm{B}} = \frac{e\hbar}{2m_{\mathrm{e}}}.\tag{1.15}
\end{equation*}
这是描述原子磁矩大小的方便单位，数值为 $9.274\times10^{-24}$~A\,m$^2$。注意式~(1.14)
中磁矩为负：由于电子带负电，其磁矩与角动量反平行。电子的旋磁比为
$\gamma=-e/2m_{\mathrm{e}}$，拉莫尔频率则为 $\omega_{\mathrm{L}}=|\gamma|B=eB/2m_{\mathrm{e}}$。

\begin{figure}
\centering
\includegraphics[width=0.42\textwidth]{fig1_05.pdf}
\caption{氢原子中绕核（单个质子）以速度 $\mathbf{v}$ 作轨道运动的电子。}
\end{figure}
\bio{Niels Bohr（1885--1962）}

\subsection{磁化强度与磁场}

磁性固体由大量带磁矩的原子组成。磁化强度（magnetization）$\mathbf{M}$ 定义为单位
体积的磁矩。通常在“连续介质近似”下考虑这个矢量：即在足够大的空间尺度上观察，
看不到单个原子磁矩带来的颗粒感。这样 $\mathbf{M}$ 可视为一个光滑的矢量场，只在
磁性固体边界处不连续。

自由空间（真空）中没有磁化。磁场可用线性相关的两个矢量场 $\mathbf{B}$ 与
$\mathbf{H}$ 描述：
\begin{equation*}
\mathbf{B} = \mu_{0}\mathbf{H},\tag{1.16}
\end{equation*}
其中 $\mu_{0}=4\pi\times10^{-7}$~H\,m$^{-1}$ 是真空磁导率。$\mathbf{B}$ 与 $\mathbf{H}$
只是彼此的标度版本：前者以特斯拉（T）度量，后者以 A\,m$^{-1}$ 度量。

在磁性固体中，$\mathbf{B}$ 与 $\mathbf{H}$ 的关系更为复杂，两个矢量场的大小和方向
都可以相差很大。一般的矢量关系是
\begin{equation*}
\mathbf{B} = \mu_{0}(\mathbf{H} + \mathbf{M}).\tag{1.17}
\end{equation*}
\mnote{由于历史原因，标准约定把 $\mathbf{B}$ 称为磁感应强度（magnetic induction）或磁通密度（magnetic flux density），把 $\mathbf{H}$ 称为磁场强度（magnetic field strength）。然而这些名称既繁琐又易误导。遵照通常的用法，我们把二者都简称为磁场，由字母 B 和 H 表明所指为何。}

当磁化强度 $\mathbf{M}$ 与磁场 $\mathbf{H}$ 呈线性关系时，该固体称为线性材料，
写作
\begin{equation*}
\mathbf{M} = \chi\mathbf{H},\tag{1.18}
\end{equation*}
$\chi$ 是无量纲量，称为磁化率（magnetic susceptibility）。此时 $\mathbf{B}$ 与
$\mathbf{H}$ 仍呈线性关系：
\begin{equation*}
\mathbf{B} = \mu_{0}(1+\chi)\mathbf{H} = \mu_{0}\mu_{\mathrm{r}}\mathbf{H},\tag{1.19}
\end{equation*}
$\mu_{\mathrm{r}}=1+\chi$ 是材料的相对磁导率。

下面是一段警示。它源于我们必须在可磁化介质中非常小心地定义场。考虑自由空间中
施加磁场 $\mathbf{B}_{\mathrm{a}}$、$\mathbf{H}_{\mathrm{a}}$ 的某个区域，二者满足
$\mathbf{B}_{\mathrm{a}}=\mu_{0}\mathbf{H}_{\mathrm{a}}$。至此一切都很简单。现在把
一块磁性固体插入该区域。固体内部的场 $\mathbf{B}_{\mathrm{i}}$、$\mathbf{H}_{\mathrm{i}}$
可能与 $\mathbf{B}_{\mathrm{a}}$、$\mathbf{H}_{\mathrm{a}}$ 大不相同，原因是固体中
所有磁矩产生的磁场。事实上 $\mathbf{B}_{\mathrm{i}}$ 与 $\mathbf{H}_{\mathrm{i}}$
都可以依赖于测量点在固体内的位置——\fn{3}{被磁化的样品也会影响其外部的磁场，正如它在内部（这里所考虑的）的影响一样；玩过条形磁铁与铁屑的人都清楚这一点。}除非是
椭球形的特殊样品（见图~1.6）。若磁场沿椭球某个主轴施加，则在样品内处处有
\begin{equation*}
\mathbf{H}_{\mathrm{i}} = \mathbf{H}_{\mathrm{a}} - N\mathbf{M},\tag{1.20}
\end{equation*}
其中 $N$ 是相应的退磁因子（demagnetizing factor，见附录 D）。为从 $\mathbf{H}_{\mathrm{a}}$
得到 $\mathbf{H}_{\mathrm{i}}$ 而需加的“修正项” $\mathbf{H}_{\mathrm{d}}=-N\mathbf{M}$
称为退磁场（demagnetizing field）。类似地
\begin{equation*}
\mathbf{B}_{\mathrm{i}} = \mu_{0}(\mathbf{H}_{\mathrm{i}} + \mathbf{M}) = \mathbf{B}_{\mathrm{a}} + \mu_{0}(1-N)\mathbf{M}.\tag{1.21}
\end{equation*}

\begin{figure}
\centering
\includegraphics[width=0.55\textwidth]{fig1_06.pdf}
\caption{主轴为 $a$、$b$、$c$ 的椭球形磁化固体样品。球（$a=b=c$）与薄板（$a,b\to\infty$，
$c=0$）是其特例。}
\end{figure}

\begin{example}{1.2}
对球形样品，$N=\frac{1}{3}$，球内场为
\begin{equation*}
\mathbf{H}_{\mathrm{i}} = \mathbf{H}_{\mathrm{a}} - \frac{\mathbf{M}}{3},\tag{1.22}
\end{equation*}
\begin{equation*}
\mathbf{B}_{\mathrm{i}} = \mathbf{B}_{\mathrm{a}} + \frac{2\mu_{0}\mathbf{M}}{3}.\tag{1.23}
\end{equation*}
当磁化强度与（插入样品前测得的）外加场 $|H_{\mathrm{a}}|=|B_{\mathrm{a}}|/\mu_0$
相比不够小时，这些退磁修正必须认真对待。不过，对弱磁性的特殊情形可以绕开这些
复杂因素：对 $\chi\ll1$ 的线性材料，有 $M\ll H$、$H_{\mathrm{i}}\approx H_{\mathrm{a}}$
及 $B_{\mathrm{i}}\approx\mu_0 H_{\mathrm{i}}$，于是可以放心地假定材料中的磁场就是
我们施加的磁场。第二、三章讨论抗磁性的较弱效应时将采用这一近似。\fn{4}{即使对这些材料，在精确的实验工作中退磁场仍然必须予以考虑。}在铁磁体中，退磁效应则总是重要的。
\end{example}

\begin{example}{1.3}
材料的内禀磁化率为
\begin{equation*}
\chi_{\text{intrinsic}} = \frac{M}{H_{\mathrm{i}}}.\tag{1.24}
\end{equation*}
但实验上测得的并不是这个内禀量：实验测的是磁化强度 $M$ 对外加场 $H_{\mathrm{a}}$
的响应，即
\begin{equation*}
\chi_{\text{experimental}} = \frac{M}{H_{\mathrm{a}}}.\tag{1.25}
\end{equation*}
两个量之间有关系
\begin{equation*}
\chi_{\text{experimental}} = \frac{M}{H_{\mathrm{i}} + NM} = \frac{M/H_{\mathrm{i}}}{1 + NM/H_{\mathrm{i}}} = \frac{\chi_{\text{intrinsic}}}{1 + N\chi_{\text{intrinsic}}}.\tag{1.26}
\end{equation*}
当 $\chi_{\text{intrinsic}}\ll1$ 时二者的区分无关紧要；当 $\chi_{\text{intrinsic}}$
接近或超过 1 时，区分可能极为重要。例如，从上方逼近居里温度的铁磁体（见第四章）
有 $\chi_{\text{intrinsic}}\to\infty$，而 $\chi_{\text{experimental}}\to1/N$。
\end{example}

此处还要提出最后一点提醒：铁磁材料可能没有净磁矩，因为它由磁畴（domains）构成。
\fn{5}{关于磁畴的更多内容见第 6.7 节。}每个畴内有均匀的磁化，但相邻畴的磁化方向
不同。因此样品看起来可能没有磁化，尽管在足够小的尺度上所有磁矩都局域地排列一致。

本章余下部分将讨论磁矩与经典力学（1.2 节）和量子力学（1.3 节）相关的若干进一步
的内容。

\section{经典力学与磁矩}

本节用纯经典的论证描述外磁场对电荷体系的影响。先考虑单个电荷，再用该结果计算
电荷体系的磁化强度。电磁学重要结果摘要见附录 B。

\subsection{正则动量}

经典力学中，电荷为 $q$、速度为 $\mathbf{v}$ 的粒子在电场 $\boldsymbol{\varepsilon}$
和磁场 $\mathbf{B}$ 中所受力为
\begin{equation*}
\mathbf{F} = q(\boldsymbol{\varepsilon} + \mathbf{v}\times\mathbf{B})\tag{1.27}
\end{equation*}
称为洛伦兹力（Lorentz force）。用这个熟悉的方程可以说明磁场如何修正带电粒子的
动量。利用 $\mathbf{F}=m\,\mathrm{d}\mathbf{v}/\mathrm{d}t$、$\mathbf{B}=\nabla\times\mathbf{A}$
和 $\boldsymbol{\varepsilon}=-\nabla V-\partial\mathbf{A}/\partial t$（$V$ 为电势，
$\mathbf{A}$ 为磁矢势，$m$ 为粒子质量），式~(1.27) 可改写为
\begin{equation*}
m\frac{\mathrm{d}\mathbf{v}}{\mathrm{d}t} = -q\nabla V - q\frac{\partial\mathbf{A}}{\partial t} + q\,\mathbf{v}\times(\nabla\times\mathbf{A}).\tag{1.28}
\end{equation*}
\mnote{矢量恒等式列表见附录 G。另注意 $\mathbf{v}$ 不随位置变化。}
利用矢量恒等式
\begin{equation*}
\mathbf{v}\times(\nabla\times\mathbf{A}) = \nabla(\mathbf{v}\cdot\mathbf{A}) - (\mathbf{v}\cdot\nabla)\mathbf{A}\tag{1.29}
\end{equation*}
化简式~(1.28)，得
\begin{equation*}
m\frac{\mathrm{d}\mathbf{v}}{\mathrm{d}t} + q\left(\frac{\partial\mathbf{A}}{\partial t} + (\mathbf{v}\cdot\nabla)\mathbf{A}\right) = -q\nabla(V - \mathbf{v}\cdot\mathbf{A}).\tag{1.30}
\end{equation*}
注意 $m\,\mathrm{d}\mathbf{v}/\mathrm{d}t$ 是在与粒子一起运动的坐标系中测得的力；
偏导数 $\partial\mathbf{A}/\partial t$ 度量空间固定点上 $\mathbf{A}$ 的变化率。
于是式~(1.30) 可写作
\begin{equation*}
\frac{\mathrm{d}}{\mathrm{d}t}(m\mathbf{v} + q\mathbf{A}) = -q\nabla(V - \mathbf{v}\cdot\mathbf{A})\tag{1.31}
\end{equation*}
其中 $\mathrm{d}\mathbf{A}/\mathrm{d}t$ 是 $\mathbf{A}$ 的对流导数（convective derivative）
\begin{equation*}
\frac{\mathrm{d}\mathbf{A}}{\mathrm{d}t} = \frac{\partial\mathbf{A}}{\partial t} + (\mathbf{v}\cdot\nabla)\mathbf{A},\tag{1.32}
\end{equation*}
它度量运动粒子所在处 $\mathbf{A}$ 的变化率。式~(1.31) 具有牛顿第二定律的形式
（即“某个像动量的量的变化率等于某个像势能的量的梯度”），由此自然地引导我们定义
正则动量（canonical momentum）
\begin{equation*}
\mathbf{p} = m\mathbf{v} + q\mathbf{A}\tag{1.33}
\end{equation*}
以及带电粒子感受到的一个依赖于速度的有效势能 $q(V-\mathbf{v}\cdot\mathbf{A})$。
无磁场（$\mathbf{A}=0$）时正则动量退化为熟悉的动量 $m\mathbf{v}$。动能仍等于
$\frac{1}{2}mv^{2}$，用正则动量可写成 $(\mathbf{p}-q\mathbf{A})^{2}/2m$。这一结果
下文就会用到，书中后面在磁场中与动能相联系的量子力学算符也写作
$(-\mathrm{i}\hbar\nabla-q\mathbf{A})^{2}/2m$。

\subsection{玻尔--范列文定理}

下一步计算固体中电子系统的净磁矩，即求磁场诱导的磁化强度。由式~(1.4)，磁化强度
正比于系统能量对外磁场的导数。\fn{6}{另见附录 E。}式~(1.27) 表明磁场对带电粒子的
力总是垂直于其速度，因此不做功，系统的能量不可能依赖于外磁场。若系统能量与外磁场
无关，则不会有磁化强度。
\bio{Niels Bohr（1885--1962）\\Hendrika J. van Leeuwen（1887--1974）\\Ludwig Boltzmann（1844--1906）}

这一思想凝结为玻尔--范列文定理（Bohr--van Leeuwen theorem）：经典体系中不存在
热平衡磁化强度。可概述其证明如下：在经典统计力学中，$N$ 个电荷均为 $q$ 的粒子
的配分函数 $Z$ 正比于
\begin{equation*}
\int\!\!\int\!\cdots\!\int \exp(-\beta E(\{\mathbf{r}_i,\mathbf{p}_i\}))\,\mathrm{d}\mathbf{r}_1\cdots\mathrm{d}\mathbf{r}_N\,\mathrm{d}\mathbf{p}_1\cdots\mathrm{d}\mathbf{p}_N,\tag{1.34}
\end{equation*}
其中 $\beta=1/k_{\mathrm{B}}T$，$k_{\mathrm{B}}$ 是玻尔兹曼常数，$T$ 是温度，
$i=1,\ldots,N$。这里 $E(\{\mathbf{r}_i,\mathbf{p}_i\})$ 是 $N$ 个带电粒子处于位置
$\mathbf{r}_1,\mathbf{r}_2,\ldots,\mathbf{r}_N$、动量 $\mathbf{p}_1,\mathbf{p}_2,\ldots,\mathbf{p}_N$
时的能量。积分于是展布在 $6N$ 维相空间（$3N$ 个位置坐标与 $3N$ 个动量坐标）上。
如上节所示，磁场的效果是把每个粒子的动量平移 $q\mathbf{A}$，因此须把 $\mathbf{p}_i$
换成 $\mathbf{p}_i-q\mathbf{A}$。动量积分限从 $-\infty$ 到 $\infty$，这个平移可以
通过移动动量积分的原点而吸收掉。所以配分函数不是磁场的函数，自由能
$F=-k_{\mathrm{B}}T\log Z$ 亦然（见附录 E）。故经典体系中磁化强度必为零。

这个结论初看令人意外。没有外磁场时电子走直线；加了磁场，路径弯曲（实为螺旋线）
并作回旋轨道（cyclotron orbits）。人们难免想说：这些方向一致的弯曲回旋轨道必然
贡献净磁矩，因而外磁场应当影响能量。但这一论证的谬误可借图~1.7 理解：图中画出
经典体系中电子在外磁场下的轨道。电子确实作回旋轨道，也确实对应净磁矩；把这些轨道
求和，得到系统边缘处净的逆时针环流（如图~1.1）。然而表面附近的电子无法完成完整的
圆环，它们与表面反复发生弹性碰撞，沿样品周界作所谓的跳跃轨道（skipping orbits）。
体内电子的逆时针电流与在表面反射、散射的电子的顺时针跳跃轨道电流恰好抵消。

因此玻尔--范列文定理似乎是对的，却与实验相抵触：许多含电子的真实体系确有净磁化
强度。那么定理的前提必定有问题——而这些前提正是经典力学！于是我们得出结论：
经典力学不足以解释磁性材料这一最基本的性质，要说明真实材料的磁性，量子理论
不可避免。下一节将较详细地讨论电子的量子力学。

\begin{figure}
\centering
\includegraphics[width=0.5\textwidth]{fig1_07.pdf}
\caption{经典体系中的电子在外磁场下在体内作回旋轨道（图中逆时针进动），贡献净的逆时针
轨道电流（见图~1.1）。这一净电流与在表面散射、沿样品顺时针进动的电子的跳跃轨道电流
恰好抵消。}
\end{figure}

\section{自旋的量子力学}

本节简要回顾与电子自旋量子力学有关的一些结果。关于角动量量子力学的更完整论述
可参见章末延伸阅读；与量子物理和原子物理相关的一些结果也列于附录 C。

\subsection{轨道角动量与自旋角动量}

1.1 节讨论的电子角动量与电子绕核的轨道运动相联系，称为轨道角动量（orbital angular
momentum）。在真实原子中它依赖于电子占据的电子态。按通常方式定义量子数 $l$ 与
$m_l$（见附录 C），沿固定轴（取 $z$ 轴）的轨道角动量分量为 $m_l\hbar$，轨道角动量
的大小\fn{7}{严格地说，well-defined（有确切定义）的量是角动量的平方和磁偶极矩的平方：算符 $\hat{\mathbf{L}}^2$ 的本征值为 $l(l+1)\hbar^2$，$\hat{L}_z$ 的本征值为 $m_l\hbar$；类似地，算符 $\hat{\boldsymbol{\mu}}^2$ 的本征值为 $l(l+1)\mu_{\mathrm{B}}^2$，$\hat{\mu}_z$ 的本征值为 $-m_l\mu_{\mathrm{B}}$。}为
$\sqrt{l(l+1)}\hbar$。于是沿 $z$ 轴的磁矩分量为 $-m_l\mu_{\mathrm{B}}$，总磁偶极矩的
大小为 $\sqrt{l(l+1)}\mu_{\mathrm{B}}$。

情况还因下述事实而更复杂：电子还具有与内禀角动量相联系的内禀磁矩。电子的内禀
角动量称为自旋（spin）。这个名字源于人们曾以为电子绕自身轴进动；但电子是点粒子，
这实在难以想象。\mnote{这一概念已经改变，但名称沿用至今。这倒也不是坏事：电子自旋的行为如此违反直觉，很难找到任何一个词能完全恰当地描述它！}

电子自旋由自旋量子数 $s$ 刻画，对电子 $s=\frac{1}{2}$。角动量任一分量只能取
$2s+1$ 个可能值：$s\hbar,(s-1)\hbar,\ldots,-s\hbar$。自旋角动量分量写作 $m_s\hbar$。
对 $s=\frac{1}{2}$ 的电子，只有两个可能值 $m_s=\pm\frac{1}{2}$，沿特定轴的角动量
分量为 $\hbar/2$ 或 $-\hbar/2$，分别称为“向上”和“向下”。电子自旋角动量的大小\fn{8}{严格地，算符 $\hat{\mathbf{S}}^2$ 的本征值是 $s(s+1)\hbar^2$。}为
$\sqrt{s(s+1)}\hbar=\sqrt{3}\hbar/2$。

自旋角动量联系着磁矩，其沿特定轴的分量为 $-g\mu_{\mathrm{B}}m_s$，大小为
$\sqrt{s(s+1)}\,g\mu_{\mathrm{B}}=\sqrt{3}\,g\mu_{\mathrm{B}}/2$。式中 $g$ 是常数，
称为 $g$ 因子（g-factor）。\mnote{$g$ 因子将在附录 C.6 中更详细地讨论。}$g$ 因子
近似为 2，因此尽管自旋是半整数的，电子内禀磁矩沿 $z$ 轴的分量\fn{9}{$\mp$ 号这样上下颠倒，是因为磁矩与角动量反平行，其根源是电子带负电：$m_s=+\frac{1}{2}$ 时磁矩为 $-\mu_{\mathrm{B}}$；$m_s=-\frac{1}{2}$ 时磁矩为 $+\mu_{\mathrm{B}}$。}为 $\mp\mu_{\mathrm{B}}$。
磁场 $\mathbf{B}$ 中电子的能量因此为
\begin{equation*}
E = g\mu_{\mathrm{B}}m_s B.\tag{1.35}
\end{equation*}
于是电子的能级在磁场中分裂 $g\mu_{\mathrm{B}}B$，称为塞曼分裂（Zeeman splitting）。
\bio{Pieter Zeeman（1865--1943）}

一般地，原子中的电子可以既有轨道角动量又有自旋角动量，二者组合起来。$g$ 因子在
真实原子中因此可取不同数值，依赖于自旋与轨道角动量的相对贡献。这一点下一章
再谈。

电子的角动量总是 $\hbar$ 的整数或半整数倍。为方便起见，在角动量算符的表达式中
略去 $\hbar$ 因子，即约定这些算符以 $\hbar$ 为单位度量角动量。本书今后定义的角动量
算符（如 $\hat{\mathbf{L}}$）都使角动量为 $\hbar\hat{\mathbf{L}}$。这可以简化后文
的表达式。
\bio{Wolfgang Pauli（1900--1958）}

\subsection{泡利自旋矩阵与旋量}

电子自旋的行为联系着一套相当奇特的代数，其基础是三个泡利自旋矩阵（Pauli spin
matrices）：
\begin{equation*}
\hat{\sigma}_x = \begin{pmatrix} 0 & 1 \\ 1 & 0 \end{pmatrix},\qquad
\hat{\sigma}_y = \begin{pmatrix} 0 & -\mathrm{i} \\ \mathrm{i} & 0 \end{pmatrix},\qquad
\hat{\sigma}_z = \begin{pmatrix} 1 & 0 \\ 0 & -1 \end{pmatrix}.\tag{1.36}
\end{equation*}
把它们看作“矩阵构成的矢量”会很方便：
\begin{equation*}
\boldsymbol{\sigma} = (\hat{\sigma}_x, \hat{\sigma}_y, \hat{\sigma}_z).\tag{1.37}
\end{equation*}

在继续之前，回顾几个可以直接代入验证的结果。设
\begin{equation*}
\mathbf{a} = \begin{pmatrix} a_1 \\ a_2 \\ a_3 \end{pmatrix}\tag{1.38}
\end{equation*}
为三分量矢量，则 $\boldsymbol{\sigma}\cdot\mathbf{a}$ 是矩阵
\begin{equation*}
\boldsymbol{\sigma}\cdot\mathbf{a} = \begin{pmatrix} a_3 & a_1-\mathrm{i}a_2 \\ a_1+\mathrm{i}a_2 & -a_3 \end{pmatrix}.\tag{1.39}
\end{equation*}
这样的矩阵可以相乘，例如
\begin{equation*}
(\boldsymbol{\sigma}\cdot\mathbf{a})(\boldsymbol{\sigma}\cdot\mathbf{b}) = \mathbf{a}\cdot\mathbf{b} + \mathrm{i}\,\boldsymbol{\sigma}\cdot(\mathbf{a}\times\mathbf{b})\tag{1.40}
\end{equation*}
以及
\begin{equation*}
(\boldsymbol{\sigma}\cdot\mathbf{a})^{2} = |\mathbf{a}|^{2}.\tag{1.41}
\end{equation*}
\mnote{注意式 1.40 与 1.41 中的每一项都是矩阵。$\mathbf{a}\cdot\mathbf{b}$ 是 $\mathbf{a}\cdot\mathbf{b}\,\mathbf{I}$ 的简写，$\mathbf{I}=\begin{pmatrix}1&0\\0&1\end{pmatrix}$ 是单位矩阵；类似地 $|\mathbf{a}|^2$ 是 $|\mathbf{a}|^2\mathbf{I}$ 的简写。}

现在定义自旋角动量算符
\begin{equation*}
\hat{\mathbf{S}} = \frac{1}{2}\hat{\boldsymbol{\sigma}}\tag{1.42}
\end{equation*}
于是
\begin{equation*}
\hat{S}_x = \frac{1}{2}\begin{pmatrix} 0 & 1 \\ 1 & 0 \end{pmatrix},\qquad
\hat{S}_y = \frac{1}{2}\begin{pmatrix} 0 & -\mathrm{i} \\ \mathrm{i} & 0 \end{pmatrix},\qquad
\hat{S}_z = \frac{1}{2}\begin{pmatrix} 1 & 0 \\ 0 & -1 \end{pmatrix}.\tag{1.43}
\end{equation*}
再次注意我们采用角动量以 $\hbar$ 为单位的约定，与电子相联系的角动量实际上是
$\hbar\hat{\mathbf{S}}$。（有些书把 $\hat{\mathbf{S}}$ 定义为 $\hbar\hat{\boldsymbol{\sigma}}/2$。）

只有 $\hat{S}_z$ 是对角的，因此当电子自旋沿 $z$ 方向时表象特别简单。$\hat{S}_z$
的本征值记为 $m_s$，取值 $m_s=\pm\frac{1}{2}$，相应本征态为 $|\!\uparrow_z\rangle$
与 $|\!\downarrow_z\rangle$：
\begin{equation*}
|\!\uparrow_z\rangle = \begin{pmatrix}1\\0\end{pmatrix}\tag{1.44}
\end{equation*}
\begin{equation*}
|\!\downarrow_z\rangle = \begin{pmatrix}0\\1\end{pmatrix}\tag{1.45}
\end{equation*}
分别对应自旋平行与反平行于 $z$ 轴。（“左右矢”记号即形如 $|\psi\rangle$ 的态写法，
见附录 C。）于是
\begin{equation*}
\hat{S}_z|\!\uparrow_z\rangle = \frac{1}{2}|\!\uparrow_z\rangle,\qquad
\hat{S}_z|\!\downarrow_z\rangle = -\frac{1}{2}|\!\downarrow_z\rangle.\tag{1.46}
\end{equation*}
自旋平行或反平行于 $x$、$y$ 轴的本征态是
\begin{equation*}
|\!\uparrow_x\rangle = \frac{1}{\sqrt{2}}\begin{pmatrix}1\\1\end{pmatrix},\qquad
|\!\downarrow_x\rangle = \frac{1}{\sqrt{2}}\begin{pmatrix}1\\-1\end{pmatrix}\tag{1.47}
\end{equation*}
\begin{equation*}
|\!\uparrow_y\rangle = \frac{1}{\sqrt{2}}\begin{pmatrix}1\\i\end{pmatrix},\qquad
|\!\downarrow_y\rangle = \frac{1}{\sqrt{2}}\begin{pmatrix}1\\-i\end{pmatrix}.\tag{1.48}
\end{equation*}
自旋波函数的这种双分量表示称为旋量（spinor）表示，这些态称为旋量。一般态可写为
\begin{equation*}
|\psi\rangle = \begin{pmatrix}a\\b\end{pmatrix} = a|\!\uparrow_z\rangle + b|\!\downarrow_z\rangle\tag{1.49}
\end{equation*}
$a$、$b$ 为复数，\fn{10}{当然，在一般情形下它们可以是位置的函数。}按惯例归一化为
\begin{equation*}
|a|^{2} + |b|^{2} = 1.\tag{1.50}
\end{equation*}

\begin{figure}
\centering
\includegraphics[width=0.5\textwidth]{fig1_08.pdf}
\caption{黎曼球面表示自旋 $\frac{1}{2}$ 粒子的自旋态。自旋矢量 $\mathbf{S}$ 位于
单位球面上。从球面南极到 $\mathbf{S}$ 的连线与水平赤道面（阴影）交于
$q=x+\mathrm{i}y$，该水平面被视为复平面。图中标出六种情形（$\mathbf{S}$ 平行或
反平行于 $x$、$y$、$z$ 轴）下复数 $q$ 的数值。}
\end{figure}
\bio{George F. B. Riemann（1826--1866）}

总自旋角动量算符 $\hat{\mathbf{S}}$ 定义为
\begin{equation*}
\hat{\mathbf{S}} = \begin{pmatrix}\hat{S}_x\\ \hat{S}_y\\ \hat{S}_z\end{pmatrix} = \mathbf{i}\hat{S}_x + \mathbf{j}\hat{S}_y + \mathbf{k}\hat{S}_z,\tag{1.51}
\end{equation*}
$\mathbf{i}$、$\mathbf{j}$、$\mathbf{k}$ 为笛卡尔单位矢量。算符 $\hat{\mathbf{S}}^{2}$
则为
\begin{equation*}
\hat{\mathbf{S}}^{2} = \hat{S}_x^{2} + \hat{S}_y^{2} + \hat{S}_z^{2}.\tag{1.52}
\end{equation*}
由于 $\hat{S}_x^2$、$\hat{S}_y^2$、$\hat{S}_z^2$ 的本征值总是
$\frac{1}{4}=(\pm\frac{1}{2})^{2}$，对任何自旋态 $|\psi\rangle$ 都有
\begin{equation*}
\hat{\mathbf{S}}^{2}|\psi\rangle = (\hat{S}_x^{2}+\hat{S}_y^{2}+\hat{S}_z^{2})|\psi\rangle = \left(\frac{1}{4}+\frac{1}{4}+\frac{1}{4}\right)|\psi\rangle = \frac{3}{4}|\psi\rangle.\tag{1.53}
\end{equation*}

其中许多结果可以推广到自旋量子数 $s>\frac{1}{2}$ 的粒子：最重要的结论是
$\hat{\mathbf{S}}^{2}$ 的本征值变为 $s(s+1)$。对本章考虑的 $s=\frac{1}{2}$，
$s(s+1)=\frac{3}{4}$，与式~(1.53) 一致。自旋算符间的对易关系为
\begin{equation*}
[\hat{S}_x, \hat{S}_y] = i\hat{S}_z\tag{1.54}
\end{equation*}
及其循环置换。用式~(1.40) 与 (1.42) 极易证明。这些算符都与 $\hat{\mathbf{S}}^{2}$
对易，故
\begin{equation*}
[\hat{\mathbf{S}}^{2}, \hat{S}_z] = 0.\tag{1.55}
\end{equation*}
因此可以同时确定总自旋与其一个分量，但不能同时确定多于一个的分量。

图~1.8 给出一个有助于思考自旋的几何构造。自旋矢量 $\mathbf{S}$ 指向三维空间中的
一点；由于量子态已归一化，$\mathbf{S}$ 位于单位球面上。从矢量 $\mathbf{S}$ 的端点
向球面南极引一条线，观察它与水平面（图~1.8 中阴影）的交点 $q$。把该水平面当作
阿甘德图（Argand diagram），$x$ 轴为实轴、$y$ 轴为虚轴，则 $q=x+\mathrm{i}y$ 是
复数。于是 $\mathbf{S}$ 的旋量表示为 $\begin{pmatrix}1\\q\end{pmatrix}$，归一化后为
\begin{equation*}
\mathbf{S} = \frac{1}{\sqrt{1+|q|^{2}}}\begin{pmatrix}1\\q\end{pmatrix}.\tag{1.56}
\end{equation*}
这个表示中的球面称为黎曼球面（Riemann sphere）。

\subsection{升降算符}

升降算符（raising and lowering operators）$\hat{S}_{+}$ 与 $\hat{S}_{-}$ 定义为
\begin{equation*}
\begin{aligned}
\hat{S}_{+} &= \hat{S}_x + i\hat{S}_y\\
\hat{S}_{-} &= \hat{S}_x - i\hat{S}_y.
\end{aligned}\tag{1.57}
\end{equation*}

算符 $\hat{A}$ 厄米（Hermitian）要求 $\hat{A}^{\dagger}=\hat{A}$，$\dagger$ 表示
伴随操作（对矩阵即“先转置再对每个元素取复共轭”）。升降算符不是厄米的（因为
$\hat{S}_{+}^{\dagger}=\hat{S}_{-}$，$\hat{S}_{-}^{\dagger}=\hat{S}_{+}$），因此不
对应可观测的量，但它们非常有用。直接运用式~(1.54) 与 (1.57) 得到对易关系：
\begin{equation*}
[\hat{S}_{+}, \hat{S}_{-}] = 2\hat{S}_z,\tag{1.58}
\end{equation*}
\begin{equation*}
[\hat{S}_z, \hat{S}_{\pm}] = \pm\hat{S}_{\pm},\tag{1.59}
\end{equation*}
以及
\begin{equation*}
[\hat{\mathbf{S}}^{2}, \hat{\mathbf{S}}_{\pm}] = 0.\tag{1.60}
\end{equation*}
另一个直接代入即可证明的有用关系是
\begin{equation*}
\hat{S}_{+}\hat{S}_{-} + \hat{S}_{-}\hat{S}_{+} = 2(\hat{S}_x^{2}+\hat{S}_y^{2})\tag{1.61}
\end{equation*}
它给出 $\hat{\mathbf{S}}^{2}$ 的一个方便表示：
\begin{align*}
\hat{\mathbf{S}}^{2} &= \hat{S}_x^{2}+\hat{S}_y^{2}+\hat{S}_z^{2}\tag{1.62}\\
&= \frac{1}{2}\left(\hat{S}_{+}\hat{S}_{-}+\hat{S}_{-}\hat{S}_{+}\right)+\hat{S}_z^{2}.\tag{1.63}
\end{align*}

写成矩阵，升降算符为
\begin{equation*}
\hat{S}_{+} = \begin{pmatrix}0 & 1\\ 0 & 0\end{pmatrix}\tag{1.64}
\end{equation*}
\begin{equation*}
\hat{S}_{-} = \begin{pmatrix}0 & 0\\ 1 & 0\end{pmatrix}.\tag{1.65}
\end{equation*}
\mnote{升降算符得名于它们对自旋态的作用。可以直接证明：$\hat{S}_+|\!\uparrow_z\rangle=0$，$\hat{S}_+|\!\downarrow_z\rangle=|\!\uparrow_z\rangle$，$\hat{S}_-|\!\uparrow_z\rangle=|\!\downarrow_z\rangle$，$\hat{S}_-|\!\downarrow_z\rangle=0$。因此升算符把自旋角动量的 $z$ 分量升高 $\hbar$，降算符把它降低 $\hbar$；若 $z$ 分量已处于最大（最小）值，$\hat{S}_+$（$\hat{S}_-$）将直接湮灭该态。}
利用式~(1.43)、(1.63)、(1.64) 与 (1.65) 即得
\begin{equation*}
\hat{\mathbf{S}}^{2} = \frac{3}{4}\begin{pmatrix}1 & 0\\ 0 & 1\end{pmatrix},\tag{1.66}
\end{equation*}
与式~(1.53) 一致。

\subsection{两个自旋的耦合}

现在考虑两个自旋 $\frac{1}{2}$ 粒子，由哈密顿量\fn{11}{式 1.67 这类相互作用在本书中将极为重要：超精细相互作用（见第二章）与海森伯交换相互作用（见第四章）都取这种形式。}
\begin{equation*}
\hat{\mathcal{H}} = A\,\hat{\mathbf{S}}^{a}\cdot\hat{\mathbf{S}}^{b},\tag{1.67}
\end{equation*}
描述的相互作用耦合，其中 $\hat{\mathbf{S}}^{a}$ 与 $\hat{\mathbf{S}}^{b}$ 是两个
粒子自旋的算符。把二者当作一个联合体时，总自旋也可用一个算符表示：
\begin{equation*}
\hat{\mathbf{S}}^{\mathrm{tot}} = \hat{\mathbf{S}}^{a}+\hat{\mathbf{S}}^{b}\tag{1.68}
\end{equation*}
于是
\begin{equation*}
(\hat{\mathbf{S}}^{\mathrm{tot}})^{2} = (\hat{\mathbf{S}}^{a})^{2}+(\hat{\mathbf{S}}^{b})^{2}+2\hat{\mathbf{S}}^{a}\cdot\hat{\mathbf{S}}^{b}.\tag{1.69}
\end{equation*}

两个自旋 $\frac{1}{2}$ 粒子的组合体的自旋量子数为 $s=0$ 或 1。
$(\hat{\mathbf{S}}^{\mathrm{tot}})^{2}$ 的本征值为 $s(s+1)$，故 $s=0$ 或 1 时分别为
0 或 2。由式~(1.53)，$(\hat{\mathbf{S}}^{a})^{2}$ 与 $(\hat{\mathbf{S}}^{b})^{2}$ 的
本征值都是 $\frac{3}{4}$。于是由式~(1.69)
\begin{equation*}
\hat{\mathbf{S}}^{a}\cdot\hat{\mathbf{S}}^{b} = \begin{cases} \frac{1}{4} & \text{若 } s=1\\[2pt] -\frac{3}{4} & \text{若 } s=0. \end{cases}\tag{1.70}
\end{equation*}
因为哈密顿量为 $\hat{\mathcal{H}}=A\hat{\mathbf{S}}^{a}\cdot\hat{\mathbf{S}}^{b}$，
系统对 $s=0$、$1$ 有两个能级：
\begin{equation*}
E = \begin{cases} \frac{A}{4} & \text{若 } s=1\\[2pt] -\frac{3A}{4} & \text{若 } s=0. \end{cases}\tag{1.71}
\end{equation*}
每个态的简并度为 $2s+1$，所以 $s=0$ 的态是单态（singlet），$s=1$ 的态是三重态
（triplet）。该态自旋的 $z$ 分量 $m_s$：单态取 0，三重态取 $-1$、$0$、$1$ 三个值
之一。

式~(1.70) 列出了 $\hat{\mathbf{S}}^{a}\cdot\hat{\mathbf{S}}^{b}$ 的本征值，描述其
本征态也有用处。先考虑如下基底：
\begin{equation*}
|\!\uparrow\uparrow\rangle,\quad |\!\uparrow\downarrow\rangle,\quad |\!\downarrow\uparrow\rangle,\quad |\!\downarrow\downarrow\rangle.\tag{1.72}
\end{equation*}
在该表示中第一个箭头指标记 $a$ 的自旋的 $z$ 分量，第二个箭头指标记 $b$ 的自旋的
$z$ 分量。$\hat{\mathbf{S}}^{a}\cdot\hat{\mathbf{S}}^{b}$ 的本征态是这些基底态的
线性组合，列于表~1.1；其计算见习题 1.9。注意 $m_s$ 等于各单个自旋 $z$ 分量之和。
而且由于本征态是原基底下态的混合，一般不可能同时知道原来两个自旋各自的 $z$ 分量
和联合体的总自旋——这是一般性的特征，在更复杂的情形中会变得更加重要。

\begin{table}
\centering
\caption{$\hat{\mathbf{S}}^{a}\cdot\hat{\mathbf{S}}^{b}$ 的本征态及相应的 $m_s$、$s$ 与 $\hat{\mathbf{S}}^{a}\cdot\hat{\mathbf{S}}^{b}$ 的本征值。}
\begin{tabular}{cccc}
\toprule
本征态 & $m_s$ & $s$ & $\hat{\mathbf{S}}^{a}\cdot\hat{\mathbf{S}}^{b}$ \\
\midrule
$|\!\uparrow\uparrow\rangle$ & $1$ & $1$ & $\frac{1}{4}$ \\[2pt]
$\dfrac{|\!\uparrow\downarrow\rangle+|\!\downarrow\uparrow\rangle}{\sqrt{2}}$ & $0$ & $1$ & $\frac{1}{4}$ \\[8pt]
$|\!\downarrow\downarrow\rangle$ & $-1$ & $1$ & $\frac{1}{4}$ \\[2pt]
$\dfrac{|\!\uparrow\downarrow\rangle-|\!\downarrow\uparrow\rangle}{\sqrt{2}}$ & $0$ & $0$ & $-\frac{3}{4}$ \\
\bottomrule
\end{tabular}
\end{table}

我们的基底式~(1.72) 从另一个角度看也不令人满意：波函数对两个电子的交换必须是
反对称的。波函数是空间函数 $\psi_{\text{space}}(\mathbf{r}_1,\mathbf{r}_2)$ 与自旋
函数 $\chi$ 的乘积，$\chi$ 是式~(1.72) 所列态的线性组合。空间波函数对电子交换可以
是对称或反对称的。例如空间波函数
\begin{equation*}
\psi_{\text{space}}(\mathbf{r}_1,\mathbf{r}_2) = \frac{\phi(\mathbf{r}_1)\xi(\mathbf{r}_2)\pm\phi(\mathbf{r}_2)\xi(\mathbf{r}_1)}{\sqrt{2}}\tag{1.73}
\end{equation*}
依 $\pm$ 号而对电子交换对称（$+$）或反对称（$-$）。这类对称性称为交换对称性
（exchange symmetry）。式~(1.73) 中 $\phi(\mathbf{r}_i)$ 与 $\xi(\mathbf{r}_i)$ 是
第 $i$ 个电子的单粒子波函数。无论空间波函数的交换对称性如何，自旋波函数 $\chi$
必须具有相反的交换对称性：空间波函数对称时 $\chi$ 必须反对称，反之亦然——这样才能
保证乘积 $\psi_{\text{space}}(\mathbf{r}_1,\mathbf{r}_2)\times\chi$ 整体反对称。

$|\!\uparrow\uparrow\rangle$ 与 $|\!\downarrow\downarrow\rangle$ 这样的态在电子
交换下显然对称；但交换 $|\!\uparrow\downarrow\rangle$ 中的两个电子得到
$|\!\downarrow\uparrow\rangle$，它并不等于 $|\!\uparrow\downarrow\rangle$ 的倍数。
因此 $|\!\uparrow\downarrow\rangle$（同理 $|\!\downarrow\uparrow\rangle$）在两电子
交换下既非对称也非反对称。所以毫不奇怪，我们需要取这两个态的线性组合作为本征态。
组合列于表~1.1：$(|\!\uparrow\downarrow\rangle+|\!\downarrow\uparrow\rangle)/\sqrt{2}$
在电子交换下对称（与其他两个 $s=1$ 态一样），而
$(|\!\uparrow\downarrow\rangle-|\!\downarrow\uparrow\rangle)/\sqrt{2}$ 是反对称的。

对交换的这种不对称性的另一后果是泡利不相容原理（Pauli exclusion principle）：
两个电子不能处于同一量子态。若两个电子处于完全相同的空间与自旋量子态（比如都处于
空间态 $\phi(\mathbf{r})$ 且都自旋向上），则其自旋波函数必对电子交换对称，空间波
函数就必为反对称，于是
\begin{equation*}
\psi_{\text{space}}(\mathbf{r}_1,\mathbf{r}_2) = \frac{\phi(\mathbf{r}_1)\phi(\mathbf{r}_2)-\phi(\mathbf{r}_2)\phi(\mathbf{r}_1)}{\sqrt{2}} = 0.\tag{1.74}
\end{equation*}
态消失——这正说明两个电子不能处于同一量子态。

\begin{figure}
\centering
\includegraphics[width=0.45\textwidth]{fig1_09.pdf}
\caption{两个电子以 $A\mathbf{S}^{a}\cdot\mathbf{S}^{b}$ 型相互作用耦合，给出三重态
（$s=1$）与单态（$s=0$）。若 $A>0$，单态在下、三重态在上；三重态可被磁场 $\mathbf{B}$
分裂为三个分量。}
\end{figure}

我们经常会遇到两个自旋通过 $A\hat{\mathbf{S}}^{a}\cdot\hat{\mathbf{S}}^{b}$ 型
相互作用耦合的情形（$A$ 为常数）。若 $A>0$，较低的能级是单态（能量 $-3A/4$），
其上是三重态（能量 $A/4$），比单态高出 $A$（见图~1.9）。磁场可把三重态分裂成
$m_s$ 不同的三个态。若 $A<0$，则三重态是最低能级。

\section*{延伸阅读}

\begin{itemize}[itemsep=0.2em]
\item B.~I.~Bleaney 与 B.~Bleaney，《Electricity and Magnetism》，OUP 1989，
对电磁学有全面的论述（另见附录 B）。
\item A.~I.~Rae，《Introduction to Quantum Mechanics》，IOP Publishing 1992，
是量子力学入门层次的清晰阐述。
\item 量子角动量的精彩论述见《Feynman Lectures in Physics》第三卷第 1--3 章，
R.~P.~Feynman，Addison-Wesley 1975。
\item J.~J.~Sakurai，《Modern Quantum Mechanics》，第 2 版，Addison-Wesley 1994，
是极好的量子力学参考书。
\end{itemize}

\section*{习题}

\exer{1.1} 计算电子（取 $g=2$）的磁矩。在磁通密度 0.3~T 的磁场中，该电子的拉莫尔
进动频率是多少？若电子自旋分别平行、反平行于磁场，能量差是多少？把这个能量差换算
成频率。

\exer{1.2} 利用式 1.43 中自旋算符的定义，证明式 1.53 及对易关系式 1.54 与 1.55。

\exer{1.3} 利用式 1.57 中升降算符的定义，证明式 1.58 与 1.61。

\exer{1.4} 利用自旋的对易关系 $[\hat{S}_x,\hat{S}_y]=\mathrm{i}\hat{S}_z$（及其循环
置换），证明
\begin{equation*}
[\hat{\mathbf{S}}\cdot\mathbf{X}, \hat{\mathbf{S}}] = \mathrm{i}\hat{\mathbf{S}}\times\mathbf{X},\tag{1.75}
\end{equation*}
其中 $\mathbf{X}$ 为矢量。

\exer{1.5} 利用式 1.58 与 1.61，证明
\begin{equation*}
\hat{S}_{\pm}|S, S_z\rangle = \sqrt{S(S+1)-S_z(S_z\pm1)}\,|S, S_z\pm1\rangle,\tag{1.76}
\end{equation*}
其中 $|S,S_z\rangle$ 表示总自旋角动量为 $S(S+1)\hbar^{2}$、自旋角动量 $z$ 分量为
$S_z\hbar$ 的态。据此证明式 1.76 的如下特例：
\begin{equation*}
\hat{S}_{-}|S, S\rangle = \sqrt{2}S\,|S, S-1\rangle\tag{1.77}
\end{equation*}
\begin{equation*}
\hat{S}_{+}|S, S-1\rangle = \sqrt{2}S\,|S, S\rangle\tag{1.78}
\end{equation*}

\exer{1.6} 若磁场 $\mathbf{B}$ 在空间中均匀，证明这与取 $\mathbf{A}=\frac{1}{2}(\mathbf{B}\times\mathbf{r})$ 相一致，并证明 $\nabla\cdot\mathbf{A}=0$。是否还有其他 $\mathbf{A}$ 的选取能给出同样的 $\mathbf{B}$？

\exer{1.7} 电子的动能算符为 $\hat{p}^{2}/2m$。用式 1.41 证明它可以改写为
\begin{equation*}
\frac{(\boldsymbol{\sigma}\cdot\hat{\mathbf{p}})^{2}}{2m_{\mathrm{e}}}.\tag{1.79}
\end{equation*}
施加磁场时须把 $\hat{\mathbf{p}}$ 换成 $\hat{\mathbf{p}}+e\mathbf{A}$。借助式 1.40
证明：把这一替换代入式 1.79 得到如下形式的动能
\begin{equation*}
\frac{(\hat{\mathbf{p}}+e\mathbf{A})^{2}}{2m_{\mathrm{e}}} + g\mu_{\mathrm{B}}\mathbf{B}\cdot\mathbf{S}\tag{1.80}
\end{equation*}
此处 $g$ 因子为 $g=2$。（注意本题中须小心运用式 1.40 与 1.41，因为 $\hat{\mathbf{p}}$
是算符，与 $\mathbf{A}$ 不对易。）

\exer{1.8} 某原子的轨道角动量为零，自旋量子数为 $\frac{1}{2}$，处于 $|\!\uparrow_z\rangle$
态。沿与 $z$ 轴成 $\theta$ 角的方向测量其角动量。证明测得角动量平行于新轴的概率为
$\cos^{2}(\theta/2)$。

\exer{1.9} 以式 1.72 的基底可以构造诸如 $\hat{\mathbf{S}}^{a}\cdot\hat{\mathbf{S}}^{b}$
等算符的矩阵表示；注意例如 $\hat{S}_z^{a}$ 这类算符只作用于波函数中与第一个自旋
有关的部分。于是
\begin{equation*}
\hat{S}_z^{a} = \frac{1}{2}\begin{pmatrix} 1 & 0 & 0 & 0\\ 0 & 1 & 0 & 0\\ 0 & 0 & -1 & 0\\ 0 & 0 & 0 & -1 \end{pmatrix}\tag{1.81}
\end{equation*}
\begin{equation*}
\hat{S}_z^{b} = \frac{1}{2}\begin{pmatrix} 1 & 0 & 0 & 0\\ 0 & -1 & 0 & 0\\ 0 & 0 & 1 & 0\\ 0 & 0 & 0 & -1 \end{pmatrix}\tag{1.82}
\end{equation*}
构造 $\hat{S}_x^{a}$、$\hat{S}_x^{b}$、$\hat{S}_y^{a}$、$\hat{S}_y^{b}$ 的类似表示，
进而证明
\begin{equation*}
\hat{\mathbf{S}}^{a}\cdot\hat{\mathbf{S}}^{b} = \frac{1}{4}\begin{pmatrix} 1 & 0 & 0 & 0\\ 0 & -1 & 2 & 0\\ 0 & 2 & -1 & 0\\ 0 & 0 & 0 & 1 \end{pmatrix}.\tag{1.83}
\end{equation*}
求该算符的本征值与本征矢，并验证与表 1.1 一致。

\exer{1.10} 对球形样品（a）水与（b）$\mathrm{MnSO_4\cdot4H_2O}$ 施加 0.5~T 的磁场。
分别求样品内部的 H 与 B 场与自由空间值偏离的比例。（水与 $\mathrm{MnSO_4\cdot4H_2O}$
的磁化率见表 2.1。）你会发现修正确实非常小。

\exer{1.11} 证明算符
\begin{equation*}
\hat{S}_{\theta,\phi} = \sin\theta\cos\phi\,\hat{S}_x + \sin\theta\sin\phi\,\hat{S}_y + \cos\theta\,\hat{S}_z,\tag{1.84}
\end{equation*}
（它表示自旋沿由球极角 $\theta$、$\phi$ 所确定方向的分量算符）本征值为
$\pm\frac{1}{2}$，本征态形如
\begin{align*}
|\!\uparrow\rangle &= \begin{pmatrix}\cos(\theta/2)\\ \sin(\theta/2)e^{i\phi}\end{pmatrix}\tag{1.85}\\
|\!\downarrow\rangle &= \begin{pmatrix}\sin(\theta/2)\\ -\cos(\theta/2)e^{i\phi}\end{pmatrix}.\tag{1.86}
\end{align*}
验证这些结果与图 1.8 的黎曼球面表示一致。再证明
\begin{equation*}
\hat{S}_{\theta,\phi}^{2} = \frac{1}{4}\begin{pmatrix}1 & 0\\ 0 & 1\end{pmatrix}.\tag{1.87}
\end{equation*}

\exer{1.12} 磁场沿 $z$ 方向时电子的哈密顿量（能量）算符为
\begin{equation*}
\hat{\mathcal{H}} = g\mu_{\mathrm{B}}\mathbf{B}\cdot\hat{\mathbf{S}} = g\mu_{\mathrm{B}}B\hat{S}_z.\tag{1.88}
\end{equation*}
含时薛定谔方程为
\begin{equation*}
\hat{\mathcal{H}}\psi(t) = \mathrm{i}\hbar\frac{\mathrm{d}\psi(t)}{\mathrm{d}t}\tag{1.89}
\end{equation*}
于是
\begin{equation*}
\psi(t) = \exp(-\mathrm{i}\hat{\mathcal{H}}t/\hbar)\,\psi(0).\tag{1.90}
\end{equation*}
利用式 1.41 证明
\begin{equation*}
\exp(\mathrm{i}\alpha\sigma_m) = I\cos\alpha + \mathrm{i}\sigma_m\sin\alpha\tag{1.91}
\end{equation*}
其中 $I$ 是单位矩阵，$\sigma_m$ 是某个泡利自旋矩阵，$\alpha$ 为实数。进而证明：
若把 $\psi(t)$ 写成旋量，
\begin{equation*}
\psi(t) = \begin{pmatrix} \exp(-\mathrm{i}g\mu_{\mathrm{B}}Bt/2\hbar) & 0\\ 0 & \exp(\mathrm{i}g\mu_{\mathrm{B}}Bt/2\hbar) \end{pmatrix}\psi(0),\tag{1.92}
\end{equation*}
并利用上一题的结果证明：自旋态如此演化，使 $\theta$ 的期望值不变而 $\phi$ 以
角频率 $geB/2m$ 旋转。这表明拉莫尔进动现象也可以从量子力学处理中导出。

\exer{1.13} 这是推导自旋进动的另一途径。从式 1.88 出发，用式 C.7 证明
\begin{align*}
\frac{\mathrm{d}}{\mathrm{d}t}\langle\hat{\mathbf{S}}\rangle &= \frac{1}{\mathrm{i}\hbar}\langle[\hat{\mathbf{S}}, \hat{\mathcal{H}}]\rangle\tag{1.93}\\
&= -\frac{g\mu_{\mathrm{B}}}{\hbar}\langle\hat{\mathbf{S}}\rangle\times\mathbf{B},\tag{1.94}
\end{align*}
它与式 1.6 类似，且
\begin{equation*}
\gamma = -\frac{g\mu_{\mathrm{B}}}{\hbar} = -\frac{ge}{2m_{\mathrm{e}}}.\tag{1.95}
\end{equation*}
负号来自电子带负电。

\exer{1.14} 本题讨论电偶极子的相应情形。（a）把电偶极矩为 $\mathbf{p}$、转动惯量
为 $I$ 的电偶极子置于电场 $\boldsymbol{\mathcal{E}}$ 中。经典地证明：$\mathbf{p}$ 与
$\boldsymbol{\mathcal{E}}$ 之间的夹角 $\theta$ 满足微分方程
\begin{equation*}
\ddot{\theta} = -\frac{p\mathcal{E}\sin\theta}{I}.\tag{1.96}
\end{equation*}
并证明 $\theta$ 很小时该方程给出简谐运动。

（b）用量子力学重做此题。考虑哈密顿量
\begin{equation*}
\hat{\mathcal{H}} = -\frac{\hbar^{2}}{2I}\frac{\partial^{2}}{\partial\theta^{2}} - p\mathcal{E}\cos\theta,\tag{1.97}
\end{equation*}
并说明为什么它是合适的哈密顿量。用式 C.7 证明
\begin{equation*}
\frac{\mathrm{d}}{\mathrm{d}t}\langle\hat{\theta}\rangle = \frac{\langle\hat{L}\rangle}{I}\tag{1.98}
\end{equation*}
其中 $\hat{L}=-\mathrm{i}\hbar\,\partial/\partial\theta$，且
\begin{equation*}
\frac{\mathrm{d}}{\mathrm{d}t}\langle\hat{L}\rangle = -p\mathcal{E}\langle\sin\theta\rangle.\tag{1.99}
\end{equation*}
进而推出
\begin{equation*}
\frac{\mathrm{d}^{2}}{\mathrm{d}t^{2}}\langle\hat{\theta}\rangle = -\frac{p\mathcal{E}}{I}\langle\sin\theta\rangle,\tag{1.100}
\end{equation*}
在适当的极限下它退化为经典表达式。把这与磁偶极子的情形作比较：二者为何不同？
为何电情形不出现自旋进动？

我们已经看到：电场中的电偶极子在电场平面内来回振荡，而磁场中的磁偶极子绕磁场
进动。对每一种情形，我们熟悉的“施加场（电场或磁场），偶极子（电的或磁的）就会
沿场排列”的观念有什么问题？
